% TOP_PROBLEM: {"id": 30006974, "title": "Irrationality of e+pi", "category_id": 1, "category": {"id": 1, "name": "number_theory", "display_name": "Number Theory", "description": "Properties of integers, prime numbers, Diophantine equations.", "slug": "number-theory", "order_index": 1, "created_at": "2026-07-31T15:26:25.670Z"}, "status": "open"}
% FIRST_POSTED: null
% ENTRY_KIND: statement_only
% Consolidated research record by GPT-6 Astra Pro, 27 September 2026.
% Mathematical corrections and supplementary checks are explicitly marked below.
% Compile with LuaLaTeX or XeLaTeX, twice. No external bibliography file is needed.
\documentclass[11pt]{article}
\usepackage[margin=25mm]{geometry}
\usepackage{fontspec}
\setmainfont{Latin Modern Roman}
\usepackage{amsmath,amssymb,mathtools}
\usepackage{unicode-math}
\setmathfont{Latin Modern Math}
\usepackage{xurl,hyperref}
\hypersetup{colorlinks=true,urlcolor=blue,
  pdftitle={Irrationality of e+pi: consolidated research attempt},
  pdfauthor={GPT-6 Astra Pro},
  pdfsubject={Conditional criteria, factorial identities, transcendence routes, and unresolved gaps}}
\setcounter{secnumdepth}{0}
\setlength{\parindent}{0pt}
\setlength{\parskip}{5pt}
\setlength{\emergencystretch}{3em}
\newcommand{\Q}{\mathbb Q}
\newcommand{\Z}{\mathbb Z}
\newcommand{\R}{\mathbb R}
\newcommand{\C}{\mathbb C}
\newcommand{\Qbar}{\overline{\mathbb Q}}
\DeclareMathOperator{\trdeg}{trdeg}
\DeclareMathOperator{\Log}{Log}
\DeclareMathOperator{\ord}{ord}
\begin{document}
% problemId: 30006974
% Canonical problem ID: 30006974
\section*{\texorpdfstring{Irrationality of e+pi}{Irrationality of e+pi}}\label{30006974}
\label{problem:pi-plus-e}

\subsection{Definitions and mathematical statement}
Let
\[
 e=\exp(1)=\sum_{k=0}^{\infty}\frac1{k!},\qquad
 \pi=4\arctan(1)=4\int_0^1\frac{dt}{1+t^2}.
\]
Here $0!=1$, and $n!=1\cdot2\cdots n$ for $n\ge1$. A real number is
\emph{rational} if it belongs to
$\Q=\{a/b:a,b\in\Z,\ b>0\}$, and is \emph{irrational} otherwise.
The exact question is to prove or disprove
\begin{equation}\label{eq:target}
 \boxed{e+\pi\notin\Q.}
\end{equation}
Equivalently, does
\[
 \forall a\in\Z\ \forall b\in\Z_{>0},\qquad b(e+\pi)\ne a
\]
hold? A disproof must supply a rational value with an exact proof of the
equality; a finite decimal approximation would not suffice.

A complex number is \emph{algebraic} if it is a zero of a nonzero polynomial
in $\Q[T]$; $\Qbar$ denotes the field of all algebraic numbers. A number
outside $\Qbar$ is \emph{transcendental}. Two numbers $x,y$ are
\emph{algebraically independent over $\Q$} if no nonzero
$P\in\Q[X,Y]$ satisfies $P(x,y)=0$.
The transcendence of $e$ and of $\pi$ separately is established background,
not a solution of \eqref{eq:target}; see
\href{https://webusers.imj-prg.fr/~michel.waldschmidt/articles/pdf/ColloquiumDeGiorgiSchanuelConjecture2014.pdf}{[E5]}.
The stronger assertions $e+\pi\notin\Qbar$, rational linear independence
of $1,e,\pi$, and algebraic independence of $e,\pi$ are kept distinct
from the selected target throughout this record.

\par\smallskip\textit{Formulation and status references:}
\href{https://arxiv.org/html/2311.00546v2}{[S1]}, Section 2.1;
\href{https://arxiv.org/abs/2606.17303}{[S2]}, abstract.
These are dated source checkpoints, not substitutes for a proof.

\subsection{Short English statement}
Can the sum of Euler's number and the circle constant be a fraction of two
integers? Establish that no such fraction exists, or exhibit one with a
rigorous proof of equality.

\subsection{Sources}
The mathematical references below were checked on 27 September 2026.
References describing a conjecture, an unresolved status, or a claimed
proof are not endorsements of a solution. The user-supplied text and the
two earlier responses are the source of the research narrative;
the external references provide the specifically attributed background.

\begin{itemize}
 \raggedright
 \interlinepenalty=10000
 \item[S1] Piotr Kowalski, \emph{Positive characteristic Ax--Schanuel},
 arXiv:2311.00546v2, revised 11 March 2025, Section 2.1.
 \url{https://arxiv.org/html/2311.00546v2}.
 PDF:
 \url{https://arxiv.org/pdf/2311.00546}.
 \textit{Explicit dated statement that irrationality of $e+\pi$ is unknown;
 also states Schanuel's conjecture.}

 \item[S2] Runlong Yu, \emph{Tail Criteria, No-Go Audits, and Ap\'ery-Type
 Certificate Obstructions for the Irrationality of $e+\pi$},
 arXiv:2606.17303, submitted 15 June 2026.
 \url{https://arxiv.org/abs/2606.17303}.
 \textit{Source retained from the first response. Abstract checked for its
 unresolved-status statement and description of factorial criteria.
 No independent validation of the preprint's entire audit is claimed.}

 \item[E1] Henry Cohn, \emph{A short proof of the simple continued fraction
 expansion of $e$}, American Mathematical Monthly 113 (2006), 57--62;
 arXiv:math/0601660v3.
 \url{https://arxiv.org/html/math/0601660v3}.
 \url{https://arxiv.org/abs/math/0601660}.
 \textit{Euler's continued fraction and the standard convergent recurrence.
 The selected-subsequence estimates and the conditional criteria are
 derived in the present record.}

 \item[E2] William Stein, \emph{Elementary Number Theory}, online course
 text, ``Convergence of Infinite Continued Fractions.''
 \url{https://wstein.org/edu/2007/spring/ent/ent-html/node65.html}.
 \textit{Continued-fraction approximation background; the error identity
 used here is also displayed explicitly in paragraph~2.}

 \item[E3] Boris Adamczewski and \'Eric Delaygue,
 \emph{Algebraic relations between sine and cosine values},
 arXiv:2406.01296v1, 3 June 2024, Theorem 3.
 \url{https://arxiv.org/html/2406.01296v1}.
 Original PDF link supplied by the user:
 \url{https://arxiv.org/pdf/2406.01296}.
 \textit{The algebraic-independence form of the Lindemann--Weierstrass
 theorem. This is applied only to algebraic exponents.}

 \item[E4] Vahagn Aslanyan, \emph{The Existential Closedness and Zilber--Pink
 Conjectures}, arXiv:2403.09304v1, 14 March 2024, Section 2.1.
 \url{https://arxiv.org/html/2403.09304v1}.
 \textit{Schanuel's conjecture and its application to the algebraic
 independence of $e$ and $\pi$ using $1$ and $i\pi$.}

 \item[E5] Michel Waldschmidt, \emph{Schanuel's Conjecture: algebraic
 independence of transcendental numbers}, Colloquium de Giorgi,
 6 May 2014, updated 20 May 2014, especially pages 1--4.
 \url{https://webusers.imj-prg.fr/~michel.waldschmidt/articles/pdf/ColloquiumDeGiorgiSchanuelConjecture2014.pdf}.
 \textit{Background on transcendence, Lindemann--Weierstrass, Schanuel's
 conjecture, and conditional consequences for familiar constants.}

 \item[E6] St\'ephane Fischler and Tanguy Rivoal,
 \emph{Relations between values of arithmetic Gevrey series, and
 applications to values of the Gamma function}, Journal of Number Theory
 261 (2024), 36--54; arXiv:2301.13518v1, Section 1.
 \url{https://arxiv.org/html/2301.13518v1}.
 Original PDF link supplied by the user:
 \url{https://arxiv.org/pdf/2301.13518}.
 \textit{Definitions of the function classes and value rings, and the
 discussion of $\mathbf E\cap\mathbf G=\Qbar$. The proposed
 $e^z$--$\arctan z$ construction is a proposed approach,
 not a theorem attributed to these authors.}

 \item[E7] Doron Zeilberger and Wadim Zudilin,
 \emph{The irrationality measure of $\pi$ is at most
 $7.103205334137\ldots$}, Moscow Journal of Combinatorics and Number Theory
 9 (2020), 407--419; arXiv:1912.06345v2.
 \url{https://arxiv.org/abs/1912.06345}.
 \textit{The stated irrationality-exponent bound. No exponential bound in
 the index of the convergents of $\pi$ is inferred from this result.}

 \item[D1] N. A. Carella, \emph{Linear Independence of Some Irrational
 Numbers}, arXiv:2007.15000v2, revision dated 27 February 2026.
 \url{https://arxiv.org/html/2007.15000v2}.
 \url{https://arxiv.org/abs/2007.15000v2}.
 \textit{Claimed proof examined, not an established result used as an
 input. The passages discussed are Section 2, Lemmas 4.2, 5.2 and 6.1,
 and Section 7. The citation refers specifically to version 2.}
\end{itemize}

\end{document}
